\documentclass[9pt,a4paper,twocolumn]{ctexart}

% --- Packages ---
\usepackage[top=1.2cm, bottom=1.2cm, left=1.2cm, right=1.2cm, columnsep=0.8cm]{geometry}
\usepackage{hyperref}
\usepackage{graphicx}
\usepackage{float}
\usepackage{longtable}
\usepackage{tabularx}
\usepackage{booktabs}
\usepackage{enumitem}
\usepackage{xcolor}
\usepackage{fancyhdr}
\usepackage{listings}
\usepackage{fontspec}
\usepackage{titlesec}
\usepackage{caption}

% --- Code block styling ---
\definecolor{codebg}{RGB}{245,245,245}
\definecolor{codeframe}{RGB}{220,220,220}
\definecolor{codekw}{RGB}{0,0,192}
\definecolor{codecomment}{RGB}{0,128,0}
\definecolor{codestring}{RGB}{163,21,21}
\definecolor{codenum}{RGB}{128,0,128}
\definecolor{codepunct}{RGB}{90,90,90}
\definecolor{badgegreen}{RGB}{34,139,34}
\definecolor{badgeblue}{RGB}{30,144,255}
\definecolor{badgeblack}{RGB}{50,50,50}

% --- Difficulty badges (rounded corners via tikz, pre-rendered in saveboxes) ---
\usepackage{tikz}
\newsavebox{\badgechubox}
\savebox{\badgechubox}{\tikz[baseline]{\node[fill=badgegreen,text=white,rounded corners=2.5pt,inner xsep=4pt,inner ysep=1.5pt,font=\sffamily\footnotesize\bfseries]{初级};}}
\newsavebox{\badgezhongbox}
\savebox{\badgezhongbox}{\tikz[baseline]{\node[fill=badgeblue,text=white,rounded corners=2.5pt,inner xsep=4pt,inner ysep=1.5pt,font=\sffamily\footnotesize\bfseries]{中级};}}
\newsavebox{\badgegaobox}
\savebox{\badgegaobox}{\tikz[baseline]{\node[fill=badgeblack,text=white,rounded corners=2.5pt,inner xsep=4pt,inner ysep=1.5pt,font=\sffamily\footnotesize\bfseries]{高级};}}
\newcommand{\lvchu}{\usebox{\badgechubox}}
\newcommand{\lvzhong}{\usebox{\badgezhongbox}}
\newcommand{\lvgao}{\usebox{\badgegaobox}}
\pdfstringdefDisableCommands{%
  \def\lvchu{[初级]}%
  \def\lvzhong{[中级]}%
  \def\lvgao{[高级]}%
}

\lstset{
  basicstyle=\ttfamily\scriptsize,
  keywordstyle=\color{codekw}\bfseries,
  commentstyle=\color{codecomment}\itshape,
  stringstyle=\color{codestring},
  backgroundcolor=\color{codebg},
  frame=single,
  rulecolor=\color{codeframe},
  breaklines=true,
  breakatwhitespace=false,
  breakindent=0pt,
  postbreak=\mbox{\textcolor{red}{$\hookrightarrow$}\space},
  numbers=left,
  numberstyle=\tiny\color{gray},
  columns=flexible,
  keepspaces=true,
  showstringspaces=false,
  tabsize=2,
  aboveskip=4pt,
  belowskip=4pt,
  extendedchars=true,
  literate={，}{{，}}1 {；}{{；}}1 {！}{{！}}1 {？}{{？}}1 {“}{{“}}1 {”}{{”}}1 {（}{{（}}1 {）}{{）}}1 {《}{{《}}1 {》}{{》}}1,
}

% --- Extra language definitions (no built-in listings support) ---
\lstdefinelanguage{json}{
  morekeywords={true,false,null},
  sensitive=true,
  morestring=[b]",
  comment=[l]{//},
  morecomment=[s]{/*}{*/},
}
\lstdefinelanguage{yaml}{
  morekeywords={true,false,null,yes,no,on,off},
  sensitive=false,
  morecomment=[l]{\#},
  morestring=[b]",
  morestring=[d]',
}
\lstdefinelanguage{http}{
  morekeywords={GET,POST,PUT,DELETE,PATCH,HEAD,OPTIONS,HTTP,Host,Accept,Authorization,Content-Type,Content-Length,User-Agent},
  sensitive=true,
  morecomment=[l]{\#},
  morestring=[b]",
}
\lstdefinelanguage{TypeScript}{
  morekeywords={abstract,any,as,async,await,break,case,catch,class,const,continue,debugger,declare,default,delete,do,else,enum,export,extends,false,finally,for,from,function,get,if,implements,import,in,instanceof,interface,keyof,let,module,namespace,never,new,null,number,of,package,private,protected,public,readonly,return,set,static,string,super,switch,symbol,this,throw,true,try,type,typeof,undefined,unknown,var,void,while,with,yield,Record,Array,Promise,AsyncIterable,ReadableStream,Date,Error,Object,JSON},
  sensitive=true,
  morecomment=[l]{//},
  morecomment=[s]{/*}{*/},
  morestring=[b]",
  morestring=[b]',
  morestring=[b]`{`},
}

% --- Compact spacing ---
\linespread{0.95}
\setlength{\parskip}{1pt plus 1pt minus 1pt}
\setlength{\parindent}{0pt}

% --- Table styling ---
\renewcommand{\arraystretch}{1.1}

% --- Heading styling (numbered) ---
\titleformat{\section}{\normalsize\bfseries}{\thesection\quad}{0em}{}
\titlespacing{\section}{0pt}{6pt}{2pt}
\titleformat{\subsection}{\small\bfseries}{\thesubsection\quad}{0em}{}
\titlespacing{\subsection}{0pt}{4pt}{1pt}
\titleformat{\subsubsection}{\small\bfseries}{\thesubsubsection\quad}{0em}{}
\titlespacing{\subsubsection}{0pt}{3pt}{1pt}

% --- Page style ---
\pagestyle{fancy}
\fancyhf{}
\fancyhead[L]{\small\emph{\leftmark}}
\fancyhead[R]{\small\thepage}
\renewcommand{\headrulewidth}{0.4pt}
\renewcommand{\sectionmark}[1]{}  % prevent sections from overwriting leftmark

% --- Hyperref setup ---
\hypersetup{
  colorlinks=true,
  linkcolor=blue,
  urlcolor=blue,
  citecolor=blue,
}

% --- Caption format ---
\DeclareCaptionFormat{myformat}{\small\bfseries #1#2#3}
\captionsetup{format=myformat}

\begin{document}

\title{2026 大模型面试题 | Agent 面试题 | RAG 面试题 | AI 应用开发面试指南（含答案与图解）}
\date{}
\maketitle
\markboth{Java 版 Agent 知识}{ Java 版 Agent 知识}

\section*{面试题目录}

{\footnotesize
\begin{tabular}{|p{\dimexpr 0.090\columnwidth-2\tabcolsep\relax}|p{\dimexpr 0.186\columnwidth-2\tabcolsep\relax}|p{\dimexpr 0.724\columnwidth-2\tabcolsep\relax}|}
\hline
\textbf{面试题模块} & \textbf{适合重点复习的人群} & \textbf{主要覆盖内容} \\
\hline
大模型基础面试题总结 & 所有准备 AI 应用开发面试的人 & Token、上下文窗口、采样参数、API 调用、流式输出、结构化输出、Function Calling、AI 应用评测 \\
\hline
AI Agent 面试题总结 & 准备 Agent、工具调用、工作流相关岗位的人 & Agent Loop、Memory、Prompt Engineering、Context Engineering、MCP、Agent Skills、Harness Engineering、Workflow、Graph、Loop \\
\hline
RAG 面试题总结 & 准备知识库问答、企业 AI 应用、搜索增强生成相关岗位的人 & RAG 基础、Embedding、向量数据库、Chunk 策略、Hybrid Search、Query Rewrite、Rerank、GraphRAG、知识库更新与评测 \\
\hline
AI 系统设计面试题总结 & 2 年以上开发者、准备社招和系统设计面试的人 & 生产级 AI 应用架构、模型网关、Prompt 管理、RAG、Memory、Tool Calling、可观测、评测、安全合规、实时语音 Agent \\
\hline
\end{tabular}
}



这 4 篇是“面试题入口”，每篇都会告诉你：


\begin{itemize}[itemsep=1pt, topsep=2pt]
  \item 这个模块的面试官到底想考什么。
  \item 高频题有哪些。
  \item 每组题背后应该掌握哪些关键点。
  \item 常见扣分点是什么。
  \item 应该回到哪篇原文继续深入学习。
\end{itemize}


建议你不要把它们当作纯题库看，而是当作“复习路线图”。题目只是入口，真正要掌握的是题目背后的工程判断。



这里说的“含答案与图解”，不是把所有内容压缩成几句标准答案，而是每篇面试题都会提供答题思路、关键点、扣分点和参考文章。更完整的图解和推导放在对应专题原文里，方便你从面试题继续深入学习。



\section*{AI 应用开发面试考什么？}


AI 应用开发面试和传统后端面试最大的区别是：它不只问你会不会调用接口，而是问你能不能把 AI 能力接入真实系统。



可以粗略分成三层。



\subsection*{第一层：大模型基础认知}


这一层是所有 AI 应用开发岗位都绕不开的基础。面试官通常会问：


\begin{itemize}[itemsep=1pt, topsep=2pt]
  \item Token 是什么？为什么中文、英文、代码消耗的 Token 不一样？
  \item 上下文窗口有什么限制？长上下文为什么不一定更好？
  \item Temperature、Top-P、Top-K 分别控制什么？生产环境怎么调？
  \item 大模型为什么会产生幻觉？有哪些工程缓解方式？
  \item JSON Mode、Structured Outputs、Function Calling 有什么区别？
\end{itemize}


这些题看起来基础，但真正要考的是工程认知。你不需要在普通应用开发面试里手推 Transformer，但必须知道这些参数会如何影响成本、延迟、稳定性、结构化输出和线上质量。



如果你发现自己只能背定义，讲不出生产里的影响，建议先看：大模型基础面试题总结。



\subsection*{第二层：AI 应用组件能力}


这一层是和“只会调 API”拉开差距的地方，主要包括 RAG、Agent、Prompt、Context、MCP、工具调用等。



高频题包括：


\begin{itemize}[itemsep=1pt, topsep=2pt]
  \item RAG 召回率低怎么排查？是 Chunk 问题、Embedding 问题，还是排序问题？
  \item Hybrid Search、Query Rewrite、Rerank 分别解决什么问题？
  \item Agent Loop 是什么？和普通工作流有什么区别？
  \item Agent Memory 怎么设计？短期记忆和长期记忆怎么区分？
  \item MCP 和 Function Calling 有什么区别？生产级 MCP Server 怎么做安全治理？
  \item Prompt Engineering 和 Context Engineering 到底差在哪？
\end{itemize}


这些题的共同点是：面试官不满足于听概念，而是会追问“你怎么落地”“出了问题怎么排查”“为什么这么选”。



如果你正在准备企业知识库、智能客服、Agent 工作流、AI 编程助手这类方向，建议重点看：


\begin{itemize}[itemsep=1pt, topsep=2pt]
  \item RAG 面试题总结
  \item AI Agent 面试题总结
\end{itemize}


\subsection*{第三层：AI 系统设计}


对于社招和有项目经验的候选人，这一层几乎必问。



面试官可能会直接给你一个开放题：


\begin{itemize}[itemsep=1pt, topsep=2pt]
  \item 如何设计一个企业级 AI 知识库问答系统？
  \item 如何设计一个生产级 Agent 平台？
  \item 如何设计一个模型网关，支持限流、熔断、降级和成本统计？
  \item 如何设计 AI 应用评测体系？Golden Set、LLM-as-Judge、Trace 回放怎么做？
  \item 如何设计一个实时语音 Agent？打断、低延迟、状态机怎么处理？
\end{itemize}


这类题考的是架构能力。你不能只说“用 LangChain 搭一个 RAG”，而要能讲清入口层、编排层、Prompt/Context、RAG、Memory、Tool、模型网关、可观测、评测、安全合规这些模块分别解决什么问题。



系统设计题建议直接看：AI 系统设计面试题总结。



\section*{怎么用这套面试题复习？}


这套面试题更适合“先建立框架，再回到原文深入”的方式。



\subsection*{1. 先用面试题建立知识地图}


先快速过一遍 4 篇面试题，不要求马上记住所有答案。第一遍的目标是知道 AI 应用开发面试会问哪些方向：


\begin{itemize}[itemsep=1pt, topsep=2pt]
  \item 大模型基础
  \item RAG
  \item Agent
  \item MCP 和工具调用
  \item Prompt 和 Context Engineering
  \item AI 系统设计
  \item AI 应用评测
  \item 实时语音 Agent
\end{itemize}


这一步能帮你避免复习时东一榔头西一棒子。



\subsection*{2. 再回到原文补底层理解}


每道题后面都贴了参考文章链接。遇到答不上来的题，不要急着背标准答案，先回到原文看完整逻辑。



比如：


\begin{itemize}[itemsep=1pt, topsep=2pt]
  \item Token、上下文窗口、采样参数不清楚，就看 《LLM 运行机制》。
  \item Function Calling、Structured Outputs、MCP 边界不清楚，就看 《大模型结构化输出详解》 和 《万字拆解 MCP 协议》。
  \item RAG 效果优化说不清楚，就看 《万字详解 RAG 检索优化》。
  \item 生产级 AI 应用架构说不清楚，就看 《AI 应用系统设计》。
\end{itemize}


面试题负责帮你定位考点，正文负责帮你补完整的因果链。



\subsection*{3. 最后用“工程表达”组织答案}


AI 面试题不要只答“是什么”，建议按这个结构组织：


\begin{enumerate}[itemsep=1pt, topsep=2pt]
  \item \textbf{先解释概念}：一句话讲清楚它是什么。
  \item \textbf{再说明问题}：它在真实系统里会带来什么影响。
  \item \textbf{接着给方案}：生产环境怎么设计、排查、优化或治理。
  \item \textbf{最后讲边界}：什么场景适用，什么场景不适用。
\end{enumerate}


比如问“RAG 召回率低怎么优化”，不要直接背 Hybrid Search、Rerank、Query Rewrite。更好的回答是：



先判断正确证据有没有进入候选池；如果没有，排查文档解析、Chunk、Embedding、Metadata、Query Rewrite；如果进入了但排得靠后，再考虑 Hybrid Search、Rerank、候选池大小和融合权重；如果证据进了上下文但答案仍然错，再看 Prompt、上下文位置、模型是否忠实使用证据和评测样本。



这类回答更像真的做过系统。



\section*{不同经验阶段怎么复习？}


先说结论：\textbf{不同经验阶段不是“看不看某个模块”的区别，而是掌握深度不同。}



即使是应届生，也建议至少了解 Agent 和 AI 系统设计的基本问题。现在很多校招项目、实习项目都会写智能客服、知识库问答、AI 助手、AI 编程工具，如果你完全不了解 Agent Loop、RAG 链路和生产级架构，面试官一追问就容易露怯。



更合理的复习方式是：所有人都要建立完整地图，只是深度分层。



\subsection*{应届生和 0-1 年}


目标不是把所有工程细节都背下来，而是能把 AI 应用开发的基本链路讲清楚。


\begin{itemize}[itemsep=1pt, topsep=2pt]
  \item 大模型基础面试题总结
  \item AI Agent 面试题总结
  \item RAG 面试题总结
  \item AI 系统设计面试题总结
\end{itemize}


这个阶段建议重点做到：


\begin{itemize}[itemsep=1pt, topsep=2pt]
  \item 大模型基础：能讲清 Token、上下文窗口、采样参数、结构化输出为什么会影响工程稳定性。
  \item RAG：能画出“文档处理 -> Chunk -> Embedding -> 向量库 -> 检索 -> 生成”的基本链路，并知道召回不准不能只改 Prompt。
  \item Agent：能说明 Agent 和普通 Chatbot、Workflow 的区别，知道 Agent Loop、Memory、Tools 是什么。
  \item 系统设计：能用简单语言描述一个 AI 知识库问答系统包含哪些模块，比如鉴权、RAG、模型调用、日志和评测。
\end{itemize}


应届生不一定要讲出复杂的模型网关、灰度回放和多 Agent 协作，但要表现出你不是只会复制 Demo，而是知道 Demo 到生产之间有工程差距。



\subsection*{2-3 年}


这个阶段要从“知道链路”升级到“能定位问题、能做取舍”。


\begin{itemize}[itemsep=1pt, topsep=2pt]
  \item 大模型基础面试题总结
  \item AI Agent 面试题总结
  \item RAG 面试题总结
  \item AI 系统设计面试题总结
\end{itemize}


这个阶段建议重点做到：


\begin{itemize}[itemsep=1pt, topsep=2pt]
  \item 大模型基础：能讲清 API 调用链路、幂等、限流、重试、结构化输出失败处理。
  \item RAG：能按文档处理、召回、排序、上下文、生成、评测这几段排查问题。
  \item Agent：能讲清 Agent Loop、Memory、MCP、Function Calling、Skills 的边界和组合方式。
  \item 系统设计：能讲一个生产级 AI 应用的核心模块，至少覆盖 Prompt 管理、RAG、Tool Calling、安全和可观测。
\end{itemize}


面试官会更关注你是否能把 AI 能力接入真实业务系统。比如“知识库更新后旧答案还在怎么办”“工具调用失败怎么降级”“如何证明新 Prompt 比旧 Prompt 更好”，这些问题要能给出工程化回答。



\subsection*{3 年以上}


这个阶段系统设计会成为重点，但大模型基础、RAG 和 Agent 仍然不能丢。区别是：你不能只讲单点技术，要能讲完整架构、治理策略和演进路线。


\begin{itemize}[itemsep=1pt, topsep=2pt]
  \item 大模型基础面试题总结
  \item AI Agent 面试题总结
  \item RAG 面试题总结
  \item AI 系统设计面试题总结
\end{itemize}


这个阶段建议重点做到：


\begin{itemize}[itemsep=1pt, topsep=2pt]
  \item 架构设计：能拆出入口层、编排层、Prompt/Context、RAG、Memory、Tool、模型网关、评测观测和安全合规模块。
  \item 治理能力：能讲清模型路由、fallback、Token 成本归因、Prompt 版本管理、权限隔离、审计日志。
  \item 质量闭环：能说明 Golden Set、Trace 回放、线上灰度、LLM-as-Judge 和人工复核怎么配合。
  \item 风险控制：能处理 Prompt 注入、工具越权、隐私泄露、RAG 权限过滤、模型供应商故障等问题。
\end{itemize}


这个阶段最容易被追问“如果上线后效果变差，你怎么定位？”“如果模型供应商限流，你怎么降级？”“如果 Agent 工具调错了怎么办？”“如何证明新 Prompt 比旧 Prompt 更好？”这些问题都需要工程闭环，而不是概念答案。



\section*{面试题与专栏的关系}

\begin{itemize}[itemsep=1pt, topsep=2pt]
  \item 面试题页是索引：定位高频考点；
  \item 专栏是正文：讲原理、工程细节和实践方案。
\end{itemize}


如果你只想临时抱佛脚，可以先刷 4 篇面试题；如果你想真正把 AI 应用开发这块补扎实，建议按专题把原文也读完。



\section*{后续会继续更新}


AI 应用开发还在快速变化，面试题也会继续更新。后面如果出现新的高频方向，比如多模态 Agent、端侧模型、AI Coding 工程化、MCP 生态实践、企业级评测平台，我也会继续补到这套面试题里。



如果你发现某个高频题还没覆盖，也欢迎在项目 issue 区留言。


\end{document}
